% =====================================================================
% Section 4: Method — EvidenceFlow
% =====================================================================
%
% Status: Draft §4 Method v7.3.2 — F_q condition aligned to code (2026-05-16)
%   v1-v5: see archive/
%   v6: science-object + true formulas (PPR on fact graph + noisy-OR)
%   v7: per-user feedback —
%        - opening compressed: one-line intro + three section descriptions
%        - §4.1 reduced to 2 points (≈20% of §4)
%        - §4.2 keeps formulas, prose tightened (≈40%)
%        - §4.3 each enhancement gets 1-2 formulas (≈40%)
%   v7.1:
%        - §4.1 explicitly couples entity extraction, facts, provenance, and graph layers
%        - §4.2 separates query-local graph G_q from local propagation graph H_q
%        - §4.3 connects expansion score to final reranking via s_base
%        - equations split to fit ACL double-column layout
%   v7.2:
%        - §4.2 explains why G_q and H_q are separated
%        - §4.2 defines local facts F_q more explicitly
%        - §4.3 defines the combined candidate set C_q and reranks over C_q
%   v7.3:
%        - §4.2 F_q condition tightened from OR to AND (couples F_q to G_q on both sides)
%        - §4.2 fact-fact link no longer says "activated" (every endpoint in G_q is)
%        - §4.3 m_q(e) given a precise set-theoretic definition with new Eq. (3)
%        - §4.3 σ clipped to non-negative via [·]_+ so expansion never subtracts
%        - §4.3 Δ_q rewritten via δ_{q,b} helper to avoid double-column align break
%   v7.3.1:
%        - §4.3 separates relation phrase r_e from embeddings \mathbf{r}_e and \mathbf{q}
%   v7.3.2:
%        - §4.2 F_q condition relaxed to provenance-only, matching local_ppr.py implementation
%          (verified against src/agsto/local_ppr.py:_fact_units_for_docs)
%        - F_q's query-locality is now justified via the upstream construction of G_q's
%          passage layer rather than a phrase-side condition
% Framing reference: research_memory/emnlp_expand_then_compose/46_*.md
% Writing discipline: research_memory/emnlp_expand_then_compose/45_*.md
%
% Method name: \methodname (defined via \newcommand in main file)
%
% =====================================================================

\section{Method: \methodname{}}
\label{sec:method}

\methodname{} is a graph-based retrieval framework for multi-hop question
answering. \S\ref{sec:method:substrate} builds a document-grounded
OpenIE graph offline, \S\ref{sec:method:retrieval} retrieves evidence by
local graph propagation, and \S\ref{sec:method:enhancement} enhances the
ranking with relation-grounded expansion and coverage-aware reranking,
producing a ranked evidence list $R_q$.

\subsection{Offline Indexing}
\label{sec:method:substrate}

Following standard OpenIE practice
\citep{angeli2015leveraging,gutierrez2025hipporag}, offline indexing
builds a document-grounded graph
$\mathcal{G}_{\mathcal{D}}=(\mathcal{V}_{\mathcal{D}},\mathcal{E}_{\mathcal{D}})$
in two steps.

\paragraph{Entity and fact extraction.}
For each document $d \in \mathcal{D}$, a language model extracts entity
mentions and OpenIE facts $f=(h,r,t)$, where $h$ and $t$ are head and tail
entities and $r$ is the relation phrase connecting them. Entity surface
forms are normalized into a phrase-node vocabulary
$\mathcal{V}_{\mathrm{phr}}$, while every fact retains its provenance
$\mathrm{src}(f)=d$. The result is a fact set $\mathcal{F}$ in which each
relational assertion remains traceable to the document from which it was
extracted.

\paragraph{Document-grounded multi-layer graph.}
$\mathcal{G}_{\mathcal{D}}$ couples two node layers. Passage nodes
$\mathcal{V}_{\mathrm{doc}}$ represent documents in $\mathcal{D}$, and
phrase nodes $\mathcal{V}_{\mathrm{phr}}$ represent extracted entities.
Three edge types connect them: passage--phrase edges attach each passage
to the entities it mentions; relational edges connect $(h,t)$ for each
fact $f=(h,r,t)$, retaining $r$ and $\mathrm{src}(f)$ as edge attributes;
and synonymy edges connect phrase nodes with highly similar dense
embeddings. Passage--phrase edges ground the entity layer in documents,
relational edges encode OpenIE assertions among entities, and provenance
links each assertion back to the passage that supports it. We pre-compute
dense embeddings for all nodes and materialize
$\mathcal{G}_{\mathcal{D}}$ as an adjacency index built once for the
corpus; query time performs no further OpenIE extraction over
$\mathcal{D}$.

\subsection{Evidence-Driven Local Graph Retrieval}
\label{sec:method:retrieval}

Given $q$, \methodname{} first induces a query-local subgraph $G_q$ from
$\mathcal{G}_{\mathcal{D}}$. It then instantiates a local propagation
graph $H_q$ over the facts and passages contained in $G_q$, and ranks
passages by personalized PageRank on $H_q$. We keep these objects
separate: $G_q$ determines which part of the offline passage--phrase
graph is considered for the query, while $H_q$ determines the state space
on which evidence mass propagates. The resulting query-conditioned mass
distribution $\pi_q$, which we call the \emph{evidence flow}, is
therefore both local to $q$ and grounded in the offline graph.

\paragraph{Query-local subgraph.}
We extract entity anchors $A(q) \subseteq \mathcal{V}_{\mathrm{phr}}$
from $q$ using a controlled language model that operates on the question
only, never on retrieved documents. A standard dense retriever provides
an entry signal $\mathrm{Entry}(q) \subseteq \mathcal{D}$. The entry
signal is used only as an activation prior: it seeds the local graph but
does not by itself determine the final ranking. The subgraph $G_q$ is
initialized with phrase nodes corresponding to $A(q)$, passages in
$\mathrm{Entry}(q)$, and any passage that is the provenance of a
relational edge incident to $A(q)$. From this seed set, $G_q$ expands
along relational, passage--phrase, and synonymy edges of
$\mathcal{G}_{\mathcal{D}}$ for a bounded number of hops. Thus $G_q$
contains the passage and phrase nodes that can participate in query-time
retrieval, while excluding the rest of the corpus graph. This keeps the
activated region compact while allowing passages reachable through
multi-hop entity chains to enter the local retrieval space even when
their surface similarity to $q$ is low.

\paragraph{Local personalized PageRank.}
From $G_q$, we construct $H_q$ with two node types: local fact nodes and
local passage nodes. A fact $f=(h,r,t) \in \mathcal{F}$ belongs to
$\mathcal{F}_q$ when its provenance $\mathrm{src}(f)$ is a passage node
in $G_q$. Because the passage layer of $G_q$ has already been activated
by the question-conditioned construction described above,
$\mathcal{F}_q$ inherits the query-locality of $G_q$ through provenance
without requiring a separate phrase-side condition. Two fact nodes are
linked when their facts share an entity endpoint, and each fact node
links to its provenance passage. This derived graph lets propagation
move through facts while still scoring passages as the retrieval output.
The seed distribution $s_q$ places mass on facts whose arguments include
an anchor in $A(q)$, weighted by entry scores of the corresponding
passages. We solve, by power iteration on $H_q$,
\begin{equation}
\label{eq:local-ppr}
\pi_q = \alpha\, s_q + (1-\alpha)\, P_q^{\!\top}\, \pi_q,
\end{equation}
where $P_q$ is the transition matrix induced by fact--fact and
fact--passage links in $H_q$. The passage score for $d$ aggregates
stationary mass on local facts grounded in $d$:
\begin{equation}
\label{eq:passage-score}
s_{\mathrm{ppr}}(d \mid q)
=
|\mathcal{F}_{q,d}|^{-1/2}
\sum_{f \in \mathcal{F}_{q,d}} \pi_q(f).
\end{equation}
Here,
$\mathcal{F}_{q,d}=\{\,f\in\mathcal{F}_q:\mathrm{src}(f)=d\,\}$.
The square-root normalization keeps the passage score robust to
deduplication of the underlying fact set. Because propagation is
restricted to $H_q$, $s_{\mathrm{ppr}}$ rewards passages connected to
$q$ through entity-relation evidence rather than passages independently
similar to $q$.

\subsection{Retrieval Enhancement with Relational Evidence}
\label{sec:method:enhancement}

Local propagation yields an initial graph-aware ranking, but two failure
modes remain. First, a passage on a multi-hop chain can receive low
$\pi_q$-mass when it is several hops away from any anchor. Second, the
top of the ranking can be dominated by passages that cluster around a
single entity neighborhood, leaving other entities or relations demanded
by $q$ under-covered. We address both with two enhancements that operate
on top of $s_{\mathrm{ppr}}$.

\paragraph{Relation-grounded expansion.}
For each high-scoring fact in $\pi_q$, we collect its incident phrase
nodes in $G_q$ and follow nearby relational edges of
$\mathcal{G}_{\mathcal{D}}$. Their provenance passages are admitted as
expansion candidates. For a relational edge $e$ between phrase endpoints
$(h_e, t_e)$ with relation phrase $r_e$ and provenance
$\mathrm{src}(e)=d$, let
\begin{equation}
\label{eq:edge-mass}
\begin{aligned}
m_q(e) = \max\bigl\{\, \pi_q(f) :{}& f=(h,r,t) \in \mathcal{F}_q, \\
& \{h, t\} \cap \{h_e, t_e\} \neq \emptyset \,\bigr\},
\end{aligned}
\end{equation}
with $m_q(e)=0$ if no such fact exists. The expansion score combines this
mass with the dense compatibility between $r_e$ and $q$:
\begin{equation}
\label{eq:expansion-score}
s_{\mathrm{exp}}(d \mid q)
=
\max_{e:\,\mathrm{src}(e)=d}
m_q(e)\,\bigl[\sigma(\mathbf{r}_e,\mathbf{q})\bigr]_{+},
\end{equation}
where $\mathbf{r}_e$ and $\mathbf{q}$ are dense embeddings of the
relation phrase on edge $e$ and the question, respectively; $\sigma$ is
cosine similarity; and $[x]_{+} = \max(x, 0)$ clips to non-negative
values so that expansion can only add evidence. We combine local-PPR and
expansion evidence into a base score
\begin{equation}
\label{eq:base-score}
s_{\mathrm{base}}(d \mid q)
=
s_{\mathrm{ppr}}(d \mid q)
+
\eta\,s_{\mathrm{exp}}(d \mid q).
\end{equation}
This step admits bridge and tail-entity passages required by the
relational structure of $q$ even when they are not lexically similar to
$q$. Let $C_q$ denote the union of passages scored by local-PPR and
passages admitted by relation-grounded expansion; passages without an
expansion path have $s_{\mathrm{exp}}(d \mid q)=0$.

\paragraph{Coverage-aware reranking.}
We extract a small set of question-side requirements
$B(q)=\{b_1,b_2,\ldots\}$ from $q$ via the same controlled language model
used for $A(q)$; each $b\in B(q)$ is an entity or relation requirement
that the evidence should help resolve. This extraction is applied to the
question only and never generates evidence chains over retrieved text.
For a passage $d$ and requirement $b$, a contribution
$\phi_{d,b}\in[0,1]$ is derived from dense similarity between the
passage and the requirement under $q$. The coverage of a selected set
$R$ for requirement $b$ is computed under a noisy-OR model,
\begin{equation}
\label{eq:noisy-or}
\mathrm{cov}_q(R \mid b)
=
1 - \prod_{d \in R}\bigl(1 - \phi_{d,b}\bigr),
\end{equation}
which saturates at $1$ once any selected passage substantially
contributes to $b$. Writing
$\delta_{q,b}(d \mid R)=\mathrm{cov}_q(R \cup \{d\}\mid b)-\mathrm{cov}_q(R\mid b)$
for the per-requirement gain, the marginal coverage gain of adding $d$ is
\begin{equation}
\label{eq:marginal-gain}
\Delta_q(d \mid R) = \sum_{b \in B(q)} \delta_{q,b}(d \mid R).
\end{equation}
We form $R_q$ greedily, at each step adding the candidate $d^*$ that
maximizes
\begin{equation}
\label{eq:greedy-rerank}
d^{*}
=
\operatorname*{arg\,max}_{d \in C_q \setminus R}
\lambda\,s_{\mathrm{base}}(d \mid q)
+
(1-\lambda)\Delta_q(d \mid R).
\end{equation}
Because the per-requirement coverage saturates, passages duplicating
already-covered neighborhoods receive diminishing marginal gain, while
passages addressing previously-uncovered requirements are promoted. The
final output is the ranked evidence list $R_q$ over passages in
$\mathcal{D}$; downstream systems may consume any prefix of $R_q$.
